\documentclass[aspectratio=169]{beamer}

% Tema y configuración
\usetheme{Madrid}
\usecolortheme{default}

% Paquetes necesarios
\usepackage[utf8]{inputenc}
\usepackage[spanish]{babel}
\usepackage{amsmath}
\usepackage{amsfonts}
\usepackage{amssymb}
\usepackage{graphicx}
\usepackage{listings}
% Habilitar UTF-8 dentro de lstlisting (acentos, ñ, etc.)
\usepackage{listingsutf8}
\usepackage{xcolor}
\usepackage{verbatim}
% Para dibujos con TikZ
\usepackage{tikz}
\usetikzlibrary{positioning,trees}

% Configuración de listings para código
\lstset{
    basicstyle=\ttfamily\footnotesize,
    keywordstyle=\color{blue}\bfseries,
    commentstyle=\color{green},
    stringstyle=\color{red},
    showstringspaces=false,
    breaklines=true,
    frame=single,
    numbers=left,
    numberstyle=\tiny\color{gray},
    inputencoding=utf8,
    % Mapeo de caracteres españoles para que no falle la compilación
    literate={á}{{\'a}}1 {é}{{\'e}}1 {í}{{\'\i}}1 {ó}{{\'o}}1 {ú}{{\'u}}1
             {Á}{{\'A}}1 {É}{{\'E}}1 {Í}{{\'I}}1 {Ó}{{\'O}}1 {Ú}{{\'U}}1
             {ñ}{{\~n}}1 {Ñ}{{\~N}}1 {ü}{{\"u}}1 {Ü}{{\"U}}1
}

% Información del documento
\title{Hoja 4: Ejercicio 19}
\subtitle{TEMA}
\author{Diego Rodríguez Cubero}
\institute{UCM}
\date{\today}

% Logo (opcional)
% \logo{\includegraphics[height=1cm]{images/logo.png}}

\begin{document}

% Diapositiva de título
\begin{frame}
    \titlepage
\end{frame}

% Tabla de contenidos
\begin{frame}{Contenidos}
    \tableofcontents
\end{frame}

% Sección 1
\section{Enunciado del problema}
\begin{frame}{Enunciado del problema}
    \begin{block}{Ejercicio 19}
        Nos dan un grafo $G = ⟨V, A⟩$ no dirigido, valorado y conexo y un ARM (árbol de recubrimiento de coste mínimo) para él. También nos dan una arista nueva $(u, v) \notin A$, con un valor asociado $c$. Se pide un algoritmo de coste en $O(|V|)$ que modifique el ARM de manera que el resultado sea un ARM para $G' = ⟨V, A \cup \{(u, v)\}⟩$. Razona la corrección y el coste.
    \end{block}
\end{frame}

% Sección 2
\section{Idea del algoritmo}
\begin{frame}{Idea del algoritmo}
    \begin{itemize}
        \item Al añadir la arista $(u, v) = e$ al ARM dado, se forma un ciclo.
        \item Para volver a tener un ARM, debemos eliminar una arista del ciclo formado, la de mayor coste.
        \item Tendremos que recorrer el ciclo para encontrar dicha arista. Esto puede hacerse encontrando el camino entre $u$ y $v$ en el ARM, y comparando su arista más cara con $e$.
        \item Si $e$ es menor que la arista más cara del ciclo, eliminamos esa arista del ARM y añadimos $e$ en su lugar.
        \item Si no, el ARM permanece igual.
    \end{itemize}
\end{frame}

% Sección 3
\section{Algoritmo}
\begin{frame}[fragile]{Algoritmo}
    \begin{lstlisting}[language = C++, caption={Actualización del ARM al añadir una arista}]
    function actualizarARM(ARM, u, v, c):
    // e = arista(u, v) con coste c
    camino = encontrarCamino(ARM, u, v) // DFS o BFS , O(V)
    aristaMax = arista de mayor coste en camino // O(V)
    if (c < coste(aristaMax)):
        eliminar aristaMax de ARM // O(1)
        añadir e a ARM // O(1)
    return ARM
    \end{lstlisting}
\end{frame}

\begin{frame}{Análisis de coste}
    \begin{itemize}
        \item Encontrar el camino entre $u$ y $v$ en el ARM puede hacerse con DFS o BFS, ambos con coste $O(V + E)$, y en un árbol $E = V - 1$, por lo que es $O(V)$.
        \item Encontrar la arista de mayor coste en el camino también cuesta $O(V)$, al tener que recorrer el camino de $V$ nodos en el peor caso.
        \item Las operaciones de eliminar y añadir aristas al ARM son $O(1)$.
        \item Por lo tanto, el coste total del algoritmo es $O(V)$.
    \end{itemize}
\end{frame}

\begin{frame}{Corrección del algoritmo}
    \begin{itemize}
        \item Al añadir la arista $(u, v)$ al ARM, se forma un ciclo.
        \item Aplicando la propiedad del corte, podemos cortar al ciclo formado en $G'=⟨V, A \cup \{(u, v)\}⟩ = G + e$, cortanto en $e$ y la arista de mayor coste del ciclo $e_o$, garantizamos que el árbol resultante sigue siendo un ARM.
        \item La propiedad del corte nos dice que exactamente una de las aristas $e, e_o$ debe pertenecer al ARM óptimo.
        \item Si $c < coste(e_o)$, entonces $e$ debe estar en el ARM óptimo, y eliminamos $e_o$.
        \item Si $c \geq coste(e_o)$, entonces $e_o$ debe estar en el ARM óptimo, y no hacemos ningún cambio.
        \item Por lo tanto, el algoritmo es correcto.
    \end{itemize}
\end{frame}

% Sección 4
\section{Conclusión}
\begin{frame}{Conclusión}
    \begin{itemize}
        \item Hemos diseñado un algoritmo eficiente para actualizar un ARM tras la adición de una nueva arista.
        \item El algoritmo tiene un coste de $O(V)$, cumpliendo con los requisitos del problema.
        \item La corrección del algoritmo se debe a la aplicación de la propiedad del corte en grafos.
    \end{itemize}
\end{frame}

\end{document}